随着城市交通规模持续增长和智慧道路建设不断推进，道路侧车辆感知已成为支撑交通组织优化、停车治理和道路空间管理的重要基础。面向长期在线部署场景，地磁传感器具有低成本、低功耗和易部署等工程优势，但其观测信息维度有限，且易受环境磁场慢变漂移、车辆速度变化、连续车流扰动和伪稳态现象影响。针对上述问题，本文面向道路侧单点三轴地磁传感器 50 Hz 连续观测数据，围绕车辆事件提取、车型分类和停车与占用状态判定展开研究，形成了由连续观测到统一事件级输入、再到上层任务识别的技术链路。

首先，针对连续三轴地磁序列中的车辆扰动事件在线提取问题，研究了由一阶差分特征提取、三轴概率融合和有限状态机约束组成的车辆事件检测方法，实现了车辆进入与离开边界的稳定提取，并形成统一的事件边界、事件片段及其持续量，为后续车型分类和停车、占用判定提供共享的事件级输入基础。

其次，针对单地磁条件下车型识别中显性信息有限以及自然车流下波形时间伸缩与局部错位并存的问题，构建了基于真实道路场景的车辆事件数据集，并从端侧轻量部署与离线增强分析两个层面开展车型三分类研究。在端侧，提取反映扰动幅值、能量、面积、持续时间及形态分布的统计特征，建立基于 Softmax 回归的轻量分类方法；在离线增强层面，引入基于动态时间规整的多模板时间对齐策略，并结合一维卷积神经网络开展增强分类。实验结果表明，端侧方案能够满足低开销部署需求，离线增强方案进一步提升了车型区分能力，尤其缓解了中型车与相邻类别之间的混淆。

最后，针对开放道路场景下环境慢漂移、连续车流和伪稳态等因素导致的停车与占用判定困难，提出了基于信号稳定点与动态参考维护的停车与占用检测方法。该方法通过窗内波动与窗间水平差的双判据提取稳定点，构建环境参考与占用参考，并结合相似性门控、状态维护与转移机制以及连续车流退化分支，实现停车确认、占用维持与释放边界的统一判定。实验结果表明，所提方法在复杂工况下仍具有较好的准确性与鲁棒性，能够有效支撑开放道路场景下的停车与占用状态检测。

总体来看，本文围绕单地磁传感器条件下的道路侧车辆感知问题，形成了由车辆事件提取到车型分类和停车与占用状态判定的连续技术链路，验证了单节点三轴地磁观测在低成本、低功耗和易部署约束下支撑道路侧多任务车辆感知的可行性，可为智慧道路精细化治理与相关系统工程部署提供方法参考。
